\PassOptionsToPackage{numbers}{natbib}
\documentclass{article}
\usepackage{iclr2024_conference}
\usepackage{luatexja-fontspec}
\usepackage{hyperref}
\usepackage{url}
\usepackage{booktabs}
\usepackage{amsmath}
\usepackage{amsfonts}
\usepackage{graphicx}
\graphicspath{{../}}

\InputIfFileExists{values.tex}{}{}
\providecommand{\airasval}[1]{\textbf{??airasval:\detokenize{#1}??}}
\providecommand{\unverified}[1]{#1}
\providecommand{\airasrecordlink}[1]{#1}

\title{SciGym Table 1 の事前登録再現（第 2 版）: 現行 API の Gemini 2.5 と GPT-4.1 を公式評価器のベンチマークパックで採点する}
\author{AIRAS}

\begin{document}
\maketitle

\begin{abstract}
SciGym \citep{duan-2025-measuring} は、反応をすべて取り除いた SBML モデルを LLM エージェントに渡し、摂動実験とコード実行を繰り返して反応機構を復元させるベンチマークである。論文の Table 1 は 6 つのフロンティアモデルを SciGym-small 137 件で各 1 回走らせた平均値を報告しているが、実行間の分散も、別環境・別時点での再現性も示していない。本研究は、公開コードの Controller をそのまま用い、現在 API で提供されている Gemini-2.5-Pro、Gemini-2.5-Flash、GPT-4.1、GPT-4.1-mini の 4 モデルで Table 1 の対応する 4 行を再現する。採点は評価層 airas-eval の SciGym ベンチマークパック \citep{airasorg-2026-airas} が行う。パックは SciGym 公式の Evaluator をコミット固定で呼び、137 件のデータを sha256 で固定し、論文の 3 指標（STE、RMS、NTS）を返す。実験前に、各モデルの平均 Simulation Trajectory Error（STE）が論文値を 0.05 より大きくは上回らないことを反証線として登録し、予測区間を $\pm 0.05$ と宣言する。Reaction Matching Score（RMS、modifier あり / なし）と Network Topology Score（NTS）は同じ実行から表として報告する。本研究は同じ設計の第 1 版（Gemini-2.5-Pro の実行途中で API 予算を超過し 25 件が失敗したまま確定した）のやり直しである。実行は RIKYU（aarch64）上で行い、論文設定からの逸脱（Gemini の GA 版、libroadrunner 2.7.0、応答トークン上限）はすべて実験前に宣言する。
\end{abstract}

\section{はじめに}
\label{sec:intro}
LLM に科学的発見のサイクル全体（実験の設計、結果の解析、機構の提案）を行わせて測るベンチマークとして、SciGym \citep{duan-2025-measuring} が提案された。BioModels 由来の SBML モデルから反応をすべて取り除いたものをエージェントに渡し、エージェントは初期濃度を変えた摂動実験を要求し、返ってきた時系列を Python で解析し、最終的な SBML を提出する。提出モデルは真のモデルと、軌道誤差（STE）と反応の一致度（RMS）で比較される。

論文の主結果である Table 1 は、3 社の pro / mini 計 6 モデルを SciGym-small の 137 件で評価した平均値である \cite[s1.p10]{duan-2025-measuring}。この評価は各モデル 1 回の実行であり、temperature 1.0 でサンプリングされる \cite[s2.p1]{h4duan-2025-scigym}。したがって Table 1 の数値がどの程度安定なのか、また同じコードを別の環境と時点で走らせて同じ数値が戻るのかは、論文からは分からない。さらに Gemini の 2 モデルはプレビュー版で評価されており \cite[s1.p1]{duan-2025-measuring}、現在は同じ版を呼び出せない。

本研究はこのギャップに対し、Table 1 のうち Gemini と OpenAI の 4 行を、公開コード \citep{h4duan-2025-scigym} の Controller をそのまま使って再現する。採点は実験コードではなく評価層 airas-eval の \texttt{scigym\_small} パック \citep{airasorg-2026-airas} が行う。パックは公式の Evaluator を含む SciGym パッケージをコミット \texttt{88a7b93} に固定し \cite[s3.p1]{airasorg-2026-airas}、137 件の入力ファイルを公式リリースの sha256 と照合し、STE・RMS・NTS を返す。第 1 版（auto-res2/scigym-table1-reproduction）では指標を実験コードが集計し、Gemini-2.5-Pro の実行途中で API ゲートウェイの予算上限に達して 25 件が失敗したまま確定した。本研究はその設計上の教訓（予算超過は run を止める、1 件ごとの強制タイムアウト、採点は評価層に任せる）を反映したやり直しである。Claude の 2 モデルは本研究が用いる API ゲートウェイで提供されていないため対象外とする。判定は AIRAS の事前登録の枠組みに従い、仮説・判定基準・予測区間・実行条件を実験前に記録へ凍結し、実験後は計測値だけを流し込む。

\section{関連研究}
\label{sec:related}
SciGym \citep{duan-2025-measuring} は、単一スキルを測る既存のベンチマークと異なり、エージェントが自ら摂動実験を設計してその結果を受け取る閉ループを評価する。評価指標は 3 つで、種と種の相互作用（同じ組は 1 回だけ数える）の F1（Network Topology Score、NTS）\cite[s1.p11]{duan-2025-measuring}、反応物と生成物の集合が一致した反応の適合率・再現率・F1（RMS。modifier まで要求する厳格版もある）\cite[s1.p8]{duan-2025-measuring}、および全種の時系列の対称平均絶対パーセント誤差（STE）\cite[s1.p9]{duan-2025-measuring} である。本研究は 3 つすべてを公式 Evaluator の実装で計算する。公式 Controller が実行時に書き出すのは STE と RMS だけなので、NTS はパックが Evaluator の関数を追加で呼んで得る。

論文は pro モデルが mini モデルを一貫して上回り \cite[s1.p6]{duan-2025-measuring}、Gemini-Pro が最良であった \cite[s1.p7]{duan-2025-measuring} と報告している。本研究はこれらの順序関係についても、判定基準には含めないが、得られた表から読み取って報告する。

\section{仮説と予測}
\label{sec:hypothesis}
以下は記録（record.json）から生成された仮説・主張・判定基準・予測区間の一覧である。実験後は Observed と Verdict が計測値から埋められる。

% AUTO-GENERATED by the update_record tool. DO NOT EDIT.
% verify_paper regenerates this file from record.json and the run outputs
% and fails on any difference, so manual edits always surface.
\noindent\textbf{H1.} SciGym 論文の Table 1 のうち Gemini-2.5-Pro、Gemini-2.5-Flash、GPT-4.1、GPT-4.1-mini の 4 行は、公開コードの Controller をそのまま用い、現在 API で提供されている同名モデル（Gemini 2 種は GA 版、GPT-4.1 2 種は論文と同じ 2025-04-14 スナップショット）で SciGym-small 137 件を最大 20 反復で解かせ、提出モデルを airas-eval の scigym\_small パック（SciGym 公式 Evaluator）で採点したとき、各モデルの平均 Simulation Trajectory Error（STE）が論文値から 0.05 以内に収まる形で再現される。
\emph{Grounded on:} \texttt{\detokenize{s1.p1}}, \texttt{\detokenize{s1.p5}}, \texttt{\detokenize{s1.p10}}.
\begin{enumerate}
\item[\textbf{C1}] 現行の Gemini-2.5-Pro（google/gemini-2.5-pro）で SciGym-small 137 件を最大 20 反復で解かせたときの平均 STE は、論文の 0.3212 を 0.05 より大きくは上回らない。
  \emph{Rationale:} Table 1 で最良とされた行がそのまま再現されることは、論文の主結果が特定のプレビュー版や実行環境に依存していないことの直接の証拠であり、h1 の Gemini-2.5-Pro 行を担う。
  \emph{Cites:} \texttt{\detokenize{s1.p5}}, \texttt{\detokenize{s1.p10}}, \texttt{\detokenize{s1.p3}}, \texttt{\detokenize{s1.p2}}, \texttt{\detokenize{s1.p7}}; criterion: \texttt{\detokenize{s1.p5}}; \texttt{\detokenize{d1}}: \texttt{\detokenize{s1.p3}}, \texttt{\detokenize{s1.p4}}, \texttt{\detokenize{s2.p1}}, \texttt{\detokenize{s3.p1}}; \texttt{\detokenize{gemini-2.5-pro}}: \texttt{\detokenize{s1.p1}}.
  \emph{Criterion:} \texttt{\detokenize{gemini-2.5-pro}}.\texttt{\detokenize{trajectory_smape}} $-$ 0.3212 $\leq 0.05$.
  \emph{Prediction:} $[-0.05, 0.05]$ (論文 Table 1 の 1 回実行の平均値。前回の再現（auto-res2/scigym-table1-reproduction）で 1 件の STE の標準偏差は約 0.3、137 件平均の標準誤差は約 0.026 だった。モデル版の差を含めても ±0.05 に収まると予測する).
  \emph{Observed:} pending.
  \emph{Verdict:} pending.
\item[\textbf{C2}] 現行の Gemini-2.5-Flash（google/gemini-2.5-flash）で SciGym-small 137 件を最大 20 反復で解かせたときの平均 STE は、論文の 0.4181 を 0.05 より大きくは上回らない。
  \emph{Rationale:} Gemini 系の mini 変種の行が再現されることで、h1 の Gemini-2.5-Flash 行を担い、c1 と合わせて Gemini 系の pro / mini の関係が論文どおりかを読めるようにする。
  \emph{Cites:} \texttt{\detokenize{s1.p5}}, \texttt{\detokenize{s1.p10}}, \texttt{\detokenize{s1.p3}}, \texttt{\detokenize{s1.p2}}, \texttt{\detokenize{s1.p6}}; criterion: \texttt{\detokenize{s1.p5}}; \texttt{\detokenize{d2}}: \texttt{\detokenize{s1.p3}}, \texttt{\detokenize{s1.p4}}, \texttt{\detokenize{s2.p1}}, \texttt{\detokenize{s3.p1}}; \texttt{\detokenize{gemini-2.5-flash}}: \texttt{\detokenize{s1.p1}}.
  \emph{Criterion:} \texttt{\detokenize{gemini-2.5-flash}}.\texttt{\detokenize{trajectory_smape}} $-$ 0.4181 $\leq 0.05$.
  \emph{Prediction:} $[-0.05, 0.05]$ (論文 Table 1 の 1 回実行の平均値。前回の再現での差は +0.028 で、137 件平均の標準誤差は約 0.028 だった).
  \emph{Observed:} pending.
  \emph{Verdict:} pending.
\item[\textbf{C3}] GPT-4.1（openai/gpt-4.1、2025-04-14 スナップショット）で SciGym-small 137 件を最大 20 反復で解かせたときの平均 STE は、論文の 0.4611 を 0.05 より大きくは上回らない。
  \emph{Rationale:} 論文と同じ日付スナップショットが使える唯一の pro モデルであり、モデル版の差を含まない最も純粋な再現条件として h1 の GPT-4.1 行を担う。
  \emph{Cites:} \texttt{\detokenize{s1.p5}}, \texttt{\detokenize{s1.p10}}, \texttt{\detokenize{s1.p3}}, \texttt{\detokenize{s1.p2}}; criterion: \texttt{\detokenize{s1.p5}}; \texttt{\detokenize{d3}}: \texttt{\detokenize{s1.p3}}, \texttt{\detokenize{s1.p4}}, \texttt{\detokenize{s2.p1}}, \texttt{\detokenize{s3.p1}}; \texttt{\detokenize{gpt-4.1}}: \texttt{\detokenize{s1.p1}}.
  \emph{Criterion:} \texttt{\detokenize{gpt-4.1}}.\texttt{\detokenize{trajectory_smape}} $-$ 0.4611 $\leq 0.05$.
  \emph{Prediction:} $[-0.05, 0.05]$ (論文 Table 1 の 1 回実行の平均値。前回の再現での差は +0.047 で、137 件平均の標準誤差は約 0.023 だった).
  \emph{Observed:} pending.
  \emph{Verdict:} pending.
\item[\textbf{C4}] GPT-4.1-mini（openai/gpt-4.1-mini、2025-04-14 スナップショット）で SciGym-small 137 件を最大 20 反復で解かせたときの平均 STE は、論文の 0.6007 を 0.05 より大きくは上回らない。
  \emph{Rationale:} 同じスナップショットが使える mini モデルの行を担い、c3 と合わせて OpenAI 系の pro / mini の関係が論文どおりかを読めるようにする。
  \emph{Cites:} \texttt{\detokenize{s1.p5}}, \texttt{\detokenize{s1.p10}}, \texttt{\detokenize{s1.p3}}, \texttt{\detokenize{s1.p2}}, \texttt{\detokenize{s1.p6}}; criterion: \texttt{\detokenize{s1.p5}}; \texttt{\detokenize{d4}}: \texttt{\detokenize{s1.p3}}, \texttt{\detokenize{s1.p4}}, \texttt{\detokenize{s2.p1}}, \texttt{\detokenize{s3.p1}}; \texttt{\detokenize{gpt-4.1-mini}}: \texttt{\detokenize{s1.p1}}.
  \emph{Criterion:} \texttt{\detokenize{gpt-4.1-mini}}.\texttt{\detokenize{trajectory_smape}} $-$ 0.6007 $\leq 0.05$.
  \emph{Prediction:} $[-0.05, 0.05]$ (論文 Table 1 の 1 回実行の平均値。前回の再現での差は +0.001 で、137 件平均の標準誤差は約 0.024 だった).
  \emph{Observed:} pending.
  \emph{Verdict:} pending.
\end{enumerate}
\noindent\emph{Assumed, so that the claims together imply H1:}
\begin{itemize}
\item 同名モデルの現行版（Gemini-2.5-Pro / Flash の GA 版、GPT-4.1 / 4.1-mini の 2025-04-14 スナップショット）が、論文で使われた版と同等の能力を持つこと（c1〜c4）
\item aarch64 向けの libroadrunner 2.7.0 と antimony 2.14.0 が、論文環境の libroadrunner 2.8.0 と同じ軌道と評価値を与えること（c1〜c4）
\item temperature 1.0 での 1 回実行の平均 STE（137 件平均）の実行間変動が 0.05 より十分小さいこと（c1〜c4）
\item Vercel AI Gateway の OpenAI 互換エンドポイント経由の各モデルが、各社の公式 API と同じ応答分布を持つこと（c1〜c4）
\item 平均 STE が Table 1 の再現性を代表すること。RMS（modifier あり / なし）と NTS は同じ実行から表として報告するが、判定基準には用いない（c1〜c4）
\end{itemize}
\noindent\textbf{Sources.}
\noindent\textbf{S1} \texttt{\detokenize{duan-2025-measuring}}: Measuring Scientific Capabilities of Language Models with a Systems Biology Dry Lab (2025).
\begin{itemize}
\item[\texttt{\detokenize{s1.p1}}] (setup) ``We evaluated six models from three frontier model families, each represented by their professional and mini variants: Anthropic’s Claude-3.7-Sonnet-20250219 (P) and Claude-3.5-Haiku-20241022 (M); Google’s Gemini-2.5-Pro-Preview-03-25 (P) and Gemini-2.5-Flash-Preview-04-17 (M); and OpenAI’s GPT-4.1-2025-04-14 (P) and GPT-4.1-Mini-2025-04-14 (M).''
\item[\texttt{\detokenize{s1.p2}}] (setup) ``Due to computational resource constraints, we restricted our evaluation to SCIGYM-small, with a total of 137 tasks.''
\item[\texttt{\detokenize{s1.p3}}] (setup) ``The agent may submit its final model and conclude the discovery process at any point, subject to a maximum iteration limit (20 in our benchmark).''
\item[\texttt{\detokenize{s1.p4}}] (setup) ``If the submitted model is invalid or cannot be simulated, we allow for 3 additional debugging iterations before evaluating against the incomplete SBML structure.''
\item[\texttt{\detokenize{s1.p5}}] (result) ``Gemini-2.5-Flash 0.4181 0.1527 0.1071 0.1217 0.2399 0.1839 0.2005 GPT-4.1-mini 0.6007 0.1516 0.1253 0.1320 0.2530 0.2313 0.2322 Claude-3.5-Haiku 0.6281 0.0858 0.0421 0.0530 0.1454 0.0805 0.0987 Gemini-2.5-Pro 0.3212 0.2138 0.1664 0.1817 0.3781 0.3219 0.3383 GPT-4.1 0.4611 0.2067 0.1597 0.1740 0.3517 0.2888 0.3038 Claude-3.7-Sonnet 0.3615 0.1780 0.1698 0.1688 0.3160 0.3170 0.3047''
\item[\texttt{\detokenize{s1.p6}}] (claim) ``Our results show that across all model families, pro variants consistently outperform their mini counterparts, with both lower STE errors and higher RMS scores.''
\item[\texttt{\detokenize{s1.p7}}] (claim) ``Among all models evaluated, Gemini-Pro achieved the best overall performance.''
\item[\texttt{\detokenize{s1.p8}}] (definition) ``We consider two reactions "matched" if they have identical sets of reactants and products. In the second, more stringent version, we additionally require matching modifiers.''
\item[\texttt{\detokenize{s1.p9}}] (definition) ``We employ the Symmetric Mean Absolute Percentage Error (SMAPE) (Makridakis, 1993), which normalizes errors relative to magnitude.''
\item[\texttt{\detokenize{s1.p10}}] (result) ``Table 1: Pro models outperform their mini counterparts in SCIGYM, with Gemini dominating in both size categories. We report the average Simulation Trajectory Error (STE) and Reaction Matching Score (RMS) across all benchmark instances.''
\item[\texttt{\detokenize{s1.p11}}] (definition) ``Network Topology Score (NTS). To assess success in recovering the correct system topology, we measure pairwise interactions between all species in the system and calculate F1 score. If identical relationships between the same species appear in multiple reactions, we count them only once.''
\end{itemize}
\noindent\textbf{S2} \texttt{\detokenize{h4duan-2025-scigym}}: h4duan/SciGym (2025).
\begin{itemize}
\item[\texttt{\detokenize{s2.p1}}] (setup) ``temperature: 1.0''
\item[\texttt{\detokenize{s2.p2}}] (setup) ``max\_iterations: 5''
\item[\texttt{\detokenize{s2.p3}}] (setup) ``eval\_debug\_rounds: 5''
\end{itemize}
\noindent\textbf{S3} \texttt{\detokenize{airasorg-2026-airas}}: airas-org/airas-eval (2026).
\begin{itemize}
\item[\texttt{\detokenize{s3.p1}}] (method) ``SCIGYM\_COMMIT = "88a7b93609e35b6ecb4eb343d816d6ff09256c6a"''
\end{itemize}


\paragraph{判定基準と予測区間の根拠.}
各主張の反証線は「（再現の平均 STE）$-$（論文値）$\le 0.05$」である。0.05 は、論文 Table 1 で隣り合うモデルの最小差（Gemini-2.5-Pro の \unverified{0.3212} と Claude-3.7-Sonnet の \unverified{0.3615} の差 \unverified{0.04}）とほぼ同じで、これより大きく外れると Table 1 の順位付けの粒度で再現に失敗したことになる。1 件の STE は $[0, 1]$ の値をとり、標準偏差を 0.25 程度と見積もると 137 件平均の標準誤差は約 0.02、2 つの 1 回実行の差の標準誤差は約 0.03 である。予測区間は $\pm 0.05$ とする。基準は悪化方向の片側なので、現行版のモデルが論文より良い結果を出しても主張は成立するが、予測区間を下に外れた場合はそのとおり報告する。

RMS の適合率・再現率・F1 も同じ実行から表として報告し、各モデルの F1 が論文値 $\pm 0.06$ に入ることを予測する（第 1 版での差は $-0.054$ から $+0.047$ であった）が、1 つの主張に置ける判定基準は 1 つであるため、判定には用いない。NTS は論文が Table 1 に載せていないため比較値が無く、表として報告するのみで予測も置かない。

\section{方法}
\label{sec:method}
\paragraph{タスク.}
各インスタンスは真の SBML（\texttt{truth.xml}）、反応を取り除いた不完全な SBML（\texttt{partial.xml}）、シミュレーション条件（\texttt{truth.sedml}）からなる。エージェントは不完全な SBML を受け取り、各ターンで実験（初期濃度の変更を指定し、真のモデルの時系列を得る）、コード実行（numpy / pandas / scipy / libsbml と、仮説モデルを走らせる \texttt{simulate} 関数が使える）、提出のいずれかを markdown 形式で返す。提出は任意の時点で行え、反復数の上限は 20 である \cite[s1.p3]{duan-2025-measuring}。提出した SBML が無効な場合は追加のデバッグ反復が許され、それでも無効なら不完全 SBML のまま採点される \cite[s1.p4]{duan-2025-measuring}。

\paragraph{実行系.}
ループ、プロンプト、実験の実行は公開コード \citep{h4duan-2025-scigym} の \texttt{Controller} をそのまま用いる。本研究が書くのは、(1) OpenAI 互換 API でモデルを呼ぶ \texttt{LLM} 実装（公式の GPT 実装と同じく、システムプロンプトと全履歴を毎回送る）、(2) インスタンスごとに 1 プロセスで Controller を起動し、137 件を並列に走らせる駆動部、(3) 各インスタンスの入力ファイルと提出 SBML を評価層の入力にまとめる部分だけである。公式コードでは 5 つの履歴がクラス属性として定義されており 1 プロセスで複数件を走らせると履歴が持ち越されるが、1 件 1 プロセスとすることでこの問題を避ける。

\paragraph{指標.}
実験コードは各インスタンスについて公式の 3 ファイル（\texttt{truth.xml}、\texttt{partial.xml}、\texttt{truth.sedml}）と提出 SBML を 1 つの入力ファイルに書くだけで、指標は計算しない。airas-eval 0.11.1 の \texttt{scigym\_small} パック \citep{airasorg-2026-airas} が、入力ファイルの内容を公式リリースの sha256 と照合したうえで、コミット \texttt{88a7b93} に固定した SciGym パッケージ \cite[s3.p1]{airasorg-2026-airas} の公式 Evaluator を呼んでシミュレーションと採点を行い、137 件の単純平均として次を返す: STE（\texttt{trajectory\_smape}）、RMS の適合率・再現率・F1（modifier なし \texttt{reaction\_*}、modifier あり \texttt{reaction\_*\_with\_modifiers}）、NTS の適合率・再現率・F1（\texttt{topology\_*}）とエッジ型別の F1。提出 SBML が無い、または無効な件は、公式 Controller と同じく不完全 SBML の採点値で平均に含める。各 run の \texttt{metrics.json} はこのレポートの指標と入力サイズ（\texttt{n\_instances}、\texttt{n\_valid\_submissions}）に、実験側が数えたトークン数を添えたものである。

\section{実験設計}
\label{sec:design}
\paragraph{モデル.}
論文の 6 モデルのうち、Gemini-2.5-Pro-Preview-03-25、Gemini-2.5-Flash-Preview-04-17、GPT-4.1-2025-04-14、GPT-4.1-Mini-2025-04-14 \cite[s1.p1]{duan-2025-measuring} に対応する現行モデルとして、Vercel AI Gateway の OpenAI 互換エンドポイント経由で \texttt{google/gemini-2.5-pro}、\texttt{google/gemini-2.5-flash}、\texttt{openai/gpt-4.1}、\texttt{openai/gpt-4.1-mini} を用いる。GPT-4.1 の 2 つは論文と同じ 2025-04-14 スナップショットである。Gemini の 2 つはプレビュー版が提供終了しているため GA 版を用いる。これが本研究の再現条件の最大の違いであり、事前に宣言する。

\paragraph{データ.}
SciGym-small の 137 件 \cite[s1.p2]{duan-2025-measuring}（Hugging Face \texttt{h4duan/scigym-sbml} の small split）を用いる。

\paragraph{設定.}
\texttt{max\_iterations} は論文どおり 20（公開 config の既定値 5 \cite[s2.p2]{h4duan-2025-scigym} は動作確認用）、\texttt{eval\_debug\_rounds} は論文本文の 3 \cite[s1.p4]{duan-2025-measuring}（公開 config の既定値は 5 \cite[s2.p3]{h4duan-2025-scigym}）、\texttt{temperature} は公開 config の 1.0 \cite[s2.p1]{h4duan-2025-scigym}、摂動は初期濃度の変更のみとする。応答のトークン上限は 32768 とする。公式の GPT 実装の既定値 8192 では、思考トークンを出すモデルが本文を返せずに失敗しうるためである。各モデル、各インスタンスを 1 回ずつ実行する。1 件 1 試行に 90 分の強制タイムアウトを設け、API エラーで終わらなかった件は最大 3 試行まで最初からやり直し、3 試行とも終わらなければ提出無しとして評価層に渡す。API ゲートウェイの予算超過（HTTP 402）を検知した場合は run 全体を失敗として止め、部分的な結果は取り込まない。

\paragraph{計算資源と環境.}
実行は Seyval を通じて RIKYU（Slurm、GB200 ノード、aarch64）で行う。GPU は使わないが、資源の単位が GPU であるため 1 GPU を要求する。環境は Dockerfile と \texttt{uv.lock} で固定する。aarch64 向けの wheel が無いため、libroadrunner は公式の 2.8.0 ではなく 2.7.0、antimony は 2.14.0 を用い、rrplugins と phrasedml は除外する。この差でシミュレーション結果がわずかに変わる可能性があり、仮説の前提として記録する。137 件は 1 件 1 プロセスで並列に実行する。

\paragraph{段階.}
sanity は 1 件（BIOMD0000000027）を 2 反復で走らせて配管を確かめる。pilot は sorted した 137 件から 14 件おきに選んだ 10 件を 20 反復で走らせ、1 件あたりの所要時間と API コストを測って全件実行の可否を判断する。主張の判定は full（137 件）の結果だけで行い、sanity と pilot の結果は記録に取り込まない。

\paragraph{run と主張の対応.}
\texttt{gemini-2.5-pro} $\to$ C1、\texttt{gemini-2.5-flash} $\to$ C2、\texttt{gpt-4.1} $\to$ C3、\texttt{gpt-4.1-mini} $\to$ C4。判定に使う指標は \texttt{trajectory\_smape} である。

% ===== airas: post-experiment sections, filled once runs exist =====

\section{結果}
\label{sec:results}
実験は未実施である。実行後にこの節へ、各 run の STE と RMS を \texttt{table1}、NTS を \texttt{topology}、完走件数とトークン消費を \texttt{runs} として、また各主張の判定基準に対する計測値を記入する。

\section{考察}
\label{sec:discussion}
実験は未実施である。

\section{結論}
\label{sec:conclusion}
実験は未実施である。

\section*{データの利用可能性}
記録は \airasrecordlink{record.json} にある。実験コード、環境の固定（Dockerfile と \texttt{uv.lock}）、各 run の結果と provenance は同じリポジトリに含まれる。

\bibliographystyle{iclr2024_conference}
\bibliography{references}

\end{document}
